\documentclass[11pt]{article}

% ======================================================================
%  Page geometry
% ======================================================================
\usepackage[a4paper,margin=1.1in]{geometry}

% ======================================================================
%  Math  (load AMS packages before the math font to avoid \Bbbk clash)
% ======================================================================
\usepackage{amsmath,amssymb,amsthm}
\usepackage{mathtools}

% ======================================================================
%  Fonts & typography  (Palatino-style text + matching math)
% ======================================================================
\usepackage[T1]{fontenc}
\usepackage{newpxtext}
\let\Bbbk\relax            % newpxmath redefines \Bbbk; avoid the clash
\usepackage{newpxmath}
\usepackage{bm}
\usepackage{microtype}

% ======================================================================
%  Tables, lists, colour, boxes, links, headings
% ======================================================================
\usepackage{booktabs}
\usepackage{enumitem}
\usepackage{graphicx}
\graphicspath{{figures/}}
\usepackage{xcolor}
\usepackage[most]{tcolorbox}
\usepackage{titlesec}
\usepackage[colorlinks=true]{hyperref}

\definecolor{accent}{HTML}{1F4E79}   % deep blue
\definecolor{accentlt}{HTML}{2E74B5} % lighter blue
\definecolor{boxbg}{HTML}{F2F6FB}    % pale blue box fill

\hypersetup{
  linkcolor=accent,
  citecolor=accent,
  urlcolor=accentlt,
}

% --- Coloured section headings ---
\titleformat{\section}
  {\normalfont\Large\bfseries\color{accent}}{\thesection}{0.7em}{}
\titleformat{\subsection}
  {\normalfont\large\bfseries\color{accentlt}}{\thesubsection}{0.6em}{}

% --- Paragraph style: open, unindented ---
\setlength{\parindent}{0pt}
\setlength{\parskip}{0.55em}

% --- Highlighted "key result" box for the central equations ---
\newtcolorbox{keybox}[1][]{
  enhanced, breakable,
  colback=boxbg, colframe=accent,
  boxrule=0.6pt, arc=3pt,
  left=10pt, right=10pt, top=8pt, bottom=8pt,
  fonttitle=\bfseries, coltitle=white,
  attach boxed title to top left={yshift=-2mm, xshift=4mm},
  boxed title style={colback=accent, arc=2pt},
  #1
}

% ======================================================================
%  Notation macros
% ======================================================================
\newcommand{\Znn}{\mathbb{Z}_{\geq 0}}            % non-negative integers
\newcommand{\Reals}{\mathbb{R}}
\newcommand{\Zsp}{\mathcal{Z}}                    % joint state space
\newcommand{\Ssp}{\mathcal{S}}                    % anatomical sites
\newcommand{\Csp}{\mathcal{C}}                    % clonotype set
\newcommand{\Rsp}{\mathcal{R}}                    % forward steps (source timepoints)
\newcommand{\Dsp}{\mathcal{D}}                    % data
\newcommand{\simplex}[1]{\Delta^{#1}}
\newcommand{\Nsrc}{N^{\mathrm{src}}}
\newcommand{\dsrc}{d^{\mathrm{src}}}
\newcommand{\E}{\mathbb{E}}
\DeclareMathOperator{\Dir}{Dirichlet}
\DeclareMathOperator{\Pois}{Poisson}

\theoremstyle{definition}
\newtheorem*{remark}{Remark}

% ======================================================================
%  Title
% ======================================================================
\title{\bfseries\color{accent}\Large A Bayesian Population-Dynamics Model for
Clonal T Cells\\[2pt] across Tissue and Phenotype in Glioblastoma}
\date{\today}

\begin{document}
\maketitle

% ======================================================================
\section{Setup}
% ======================================================================

We observe annotated T cells from multiple patients, anatomical sites, and
study timepoints. Each observed cell carries five labels: a patient, a TCR
clonotype, a study timepoint, an anatomical site, and a discrete T-cell
phenotype state.

The data are represented as aggregate
clone-by-timepoint-by-site-by-phenotype counts, from which we construct
clone-level observations across forward study steps.

\subsection{Indices}

We use the following indices throughout.
\begin{itemize}[leftmargin=1.4em, itemsep=2pt]
  \item Patients: $i = 1,\ldots,I$, with $I = 6$.
  \item TCR clonotypes within patient $i$: $c = 1,\ldots,C_i$.
  \item Shared study timepoints: $r \in \{1,\ldots,R\}$, ordered as
        $\tau_1 < \tau_2 < \cdots < \tau_R$. The timepoint labels are
        comparable across patients. Currently $R = 6$; the cohort is being
        extended with additional timepoints.
  \item Anatomical site: $s \in \Ssp$, with
        $\Ssp = \{\mathrm{PBMC},\, \mathrm{CSF},\, \mathrm{Tumor}\}$; write
        $S = |\Ssp| = 3$.
  \item Annotated T-cell phenotype state: $k \in \{1,\ldots,K\}$. In the
        current application $K = 13$, and all phenotypes share a single
        phenotype space.
\end{itemize}

\subsection{Observed count tensor}

For each patient $i$, clonotype $c$, timepoint $r$, site $s$, and phenotype
$k$, let
\[
  Y_{icrsk} \in \Znn
\]
be the observed number of cells assigned to that combination: how many cells
from patient $i$, clonotype $c$, and timepoint $r$ were sampled from site $s$
and assigned phenotype $k$. Collecting the site and phenotype axes,
\[
  Y_{icr} \in \Znn^{\,|\Ssp| \times K}
\]
is the tissue-by-phenotype count tensor for clonotype $c$ in patient $i$ at
timepoint $r$.

Each tissue sample also carries a sequencing depth and a missingness flag. For
patient $i$, timepoint $r$, and tissue $s$, let
\[
  d_{i,r,s} \in \Znn
\]
be the per-sample sequencing depth: the number of TCR$^+$ T cells in that
$(\text{patient},\text{timepoint},\text{tissue})$ sample. It is the exposure used below. Let
\[
  m_{i,r,s} \in \{0,1\}
\]
be the missingness indicator, equal to $1$ when that sample was profiled and
$0$ otherwise.

\subsection{Joint tissue--phenotype state space}

The model operates on the joint tissue--phenotype state space
\[
  \Zsp = \Ssp \times \{1,\ldots,K\},
  \qquad
  z = (s,k),
\]
where $s$ is the tissue and $k$ is the phenotype. Since $|\Ssp| = 3$ and
$K = 13$, the number of joint states is
\[
  L \;=\; |\Zsp| \;=\; |\Ssp|\,K \;=\; 39 .
\]
The model therefore tracks clone counts across a $39$-state
tissue--phenotype space. When describing transitions we write the source state
as $z = (a,u)$ and the destination state as $z' = (b,v)$, where $a,b \in \Ssp$
are source and destination tissues and $u,v \in \{1,\ldots,K\}$ are source and
destination phenotypes.

\subsection{Valid forward steps}

Forward steps are defined per patient. For patient $i$, let
$\Rsp_i \subseteq \{1,\ldots,R-1\}$ be the set of source timepoints at which
patient $i$ has a valid forward step. We index a step by its source timepoint
$r \in \Rsp_i$; the step runs from timepoint $r$ to the next sampled timepoint,
written $r+1$. Thus $r$ is not a matrix index---it simply labels a forward step
within patient $i$ by where it starts.

\begin{remark}[Skipped timepoints]
Writing the destination as $r+1$ assumes consecutive sampling. When a patient is
missing a timepoint---so that, say, $T_1$ and $T_3$ are paired---$r+1$ should be
read as the \emph{next available} timepoint for that patient. Which timepoints
are paired is a data-preprocessing detail and does not change anything below.
\end{remark}

Each $r \in \Rsp_i$ is one observed forward step. Steps such as
$T_1 \!\to\! T_2$, $T_2 \!\to\! T_3$, and $T_3 \!\to\! T_4$ are treated as
repeated observations of the \emph{same} forward process: the mean matrix is
shared across all patients and steps, and the indices $i$ and $r$ merely
organize observations rather than selecting a separate matrix.

\subsection{Clone-level observations within a step}

Fix a patient $i$, a valid step $r \in \Rsp_i$, and a clonotype $c$ observed for
that patient. The source and destination count vectors over the joint state
space are read directly off the count tensor at timepoints $r$ and $r+1$,
\[
  x_{irc}(s,k) = Y_{i\, c\, r\, s\, k},
  \qquad
  y_{irc}(s,k) = Y_{i\, c\, (r+1)\, s\, k},
  \qquad (s,k)\in\Zsp,
\]
so $x_{irc} \in \Znn^{L}$ and $y_{irc} \in \Znn^{L}$.

The source enters the model as its depth-normalized composition. The
\emph{depth-rescaled source}
\[
  \tilde{x}_{irc}(z) = \frac{x_{irc}(z)}{\dsrc_{ir,s}},
  \qquad s = \mathrm{tissue}(z),
  \qquad \dsrc_{ir,s} = d_{i,r,s},
\]
is the fraction of the source sample's TCR$^+$ T cells in bin $z$ (tissue
$s = \mathrm{tissue}(z)$). The per-sample depth $d$ is retained rather than
normalized away, so that an observed destination zero can be read as contraction
or clonal loss where depth was high versus undersampling where depth was
low---the separation the depth offset provides.

We keep every clonotype with at least one cell at the source timepoint $r$,
\[
  \Csp_{ir}
  =
  \bigl\{\,
    c \in \{1,\ldots,C_{i}\}
    \;:\;
    \Nsrc_{irc} > 0
  \,\bigr\},
  \qquad
  \Nsrc_{irc} = \sum_{z\in\Zsp} x_{irc}(z) .
\]
A clone that is absent at the destination timepoint is kept: its absence is
informative---contraction or clonal loss---not a reason to drop it. This gives
the data hierarchy
\[
  i
  \;\longrightarrow\;
  r \in \Rsp_i
  \;\longrightarrow\;
  c \in \Csp_{ir},
\]
patients contain valid forward steps, and each forward step has a retained set
of clonotypes present at the source timepoint. A \emph{clone-level forward
observation} is the triple $(i,r,c)$; for compact notation we write
\[
  j \;\equiv\; (i,r,c),
  \qquad
  x_j = x_{irc},
  \qquad
  \tilde{x}_j = \tilde{x}_{irc},
  \qquad
  y_j = y_{irc} .
\]

Each destination state $z'$ carries the destination depth and missingness flag
of its tissue $s' = \mathrm{tissue}(z')$:
$d_{j,s'} = d_{i,\,r+1,\,s'}$ and $m_{j,s'} = m_{i,\,r+1,\,s'}$.

\subsection{Two kinds of destination zero}

Destination zeros are part of the likelihood specification, and the two kinds
are treated differently.
\begin{enumerate}[leftmargin=2em, itemsep=3pt]
  \item \textbf{Missing sample} ($m_{j,s'} = 0$). The destination states in
        tissue $s'$ were not profiled at the destination timepoint. They are
        excluded from the likelihood entirely---masked out, never entered as a
        zero count.
  \item \textbf{Depth-zero} ($m_{j,s'} = 1$, clone not captured). A genuine
        observed zero, modeled natively by the Negative-Binomial likelihood below: a small
        mean $d_{j,s'}\,\mu_j(z')$ emits zero with positive probability while
        $\mu_j(z') > 0$. No zero-inflation component is added.
\end{enumerate}
The depth-zero mechanism is what permits a clone to be absent at one step and
present at the next.

\subsection{Setup summary}

The setup produces the clone-level forward observations
\[
  \Dsp = \bigl\{ (\tilde{x}_j,\, y_j,\, d_j,\, m_j,\, i_j,\, r_j,\, c_j) \bigr\}_{j=1}^{J},
  \qquad j \equiv (i_j, r_j, c_j),
\]
each carrying a depth-rescaled source vector, a destination count vector, the
per-tissue destination depths $d_j = (d_{j,s'})_{s'\in\Ssp}$ and missingness
flags $m_j = (m_{j,s'})_{s'\in\Ssp}$, and a patient / step / clonotype identity.
The count vectors live in the joint tissue--phenotype state space
$\Zsp = \Ssp \times \{1,\ldots,K\}$, with $|\Zsp| = 39$.

\subsection{Notation summary}

Table~\ref{tab:notation} collects the core modeling symbols.

\begin{table}[h]
\centering
\renewcommand{\arraystretch}{1.25}
\begin{tabular}{@{}ll@{}}
\toprule
\textbf{Symbol} & \textbf{Meaning}\\
\midrule
$i = 1,\ldots,I$              & patient index\\
$c = 1,\ldots,C_i$           & TCR clonotype within patient $i$\\
$r \in \{1,\ldots,R\}$       & shared study timepoint, $\tau_1<\cdots<\tau_R$\\
$s \in \Ssp$                 & anatomical site, $\Ssp=\{\mathrm{PBMC},\mathrm{CSF},\mathrm{Tumor}\}$\\
$k \in \{1,\ldots,K\}$       & phenotype state ($K=13$)\\
$z = (s,k) \in \Zsp$         & joint tissue--phenotype state ($L=|\Zsp|=39$)\\
$Y_{icrsk}$                  & observed count tensor\\
$d_{i,r,s}$                  & per-sample sequencing depth (count exposure)\\
$m_{i,r,s} \in \{0,1\}$      & sample missingness indicator\\
$\Rsp_i$                     & valid forward steps for patient $i$ ($\subseteq\{1,\ldots,R-1\}$)\\
$r \in \Rsp_i$               & forward step, from source timepoint $r$ to $r+1$\\
$\Csp_{ir}$                  & clonotypes present at the source timepoint of step $r$\\
$j = (i_j,r_j,c_j)$          & clone-level observation, $j=1,\ldots,J$\\
$x_j,\,y_j \in \Znn^{L}$     & source / destination count vectors\\
$\tilde{x}_j$                & depth-rescaled source, $\tilde{x}_j(z)=x_j(z)/\dsrc_{ir,s}$\\
$g_{a,u} \geq 0$             & per-source-state expansion factor (see Rmk.~\ref{rmk:anchor})\\
$\pi_{(a,u)} \in \simplex{S-1}$ & destination-tissue distribution (trafficking)\\
$\Phi_{(a,u),b} \in \simplex{K-1}$ & destination-phenotype distribution per route (switching)\\
$M_{(a,u),(b,v)}=g\,\pi\,\Phi$ & assembled mean matrix $=\operatorname{diag}(g)T$ (free row sums)\\
$T_{(a,u),(b,v)}=\pi\,\Phi$   & transition matrix, row-stochastic\\
$g_z = \sum_{z'} M_{zz'}$    & row sum of $M$, equal to the expansion factor $g_{a,u}$\\
$\mu_j = \tilde{x}_j M \in \Reals_{\geq 0}^{L}$ & destination intensity vector\\
$\phi=\exp(\log\phi)$        & Negative-Binomial concentration\\
$(\mu_g,\sigma_g,\alpha_a,\beta,\sigma_\phi)$ & prior hyperparameters (log-normal / Dirichlet / NB)\\
$w_{jzz'}$                   & backward attribution probability\\
$\rho^a(u)$                  & empirical source-cell phenotype distribution in tissue $a$ (pooled)\\
$\theta(z)$                  & clone source composition, source cells as a distribution over states\\
$\Dsp$                       & full dataset\\
\bottomrule
\end{tabular}
\caption{Core modeling notation.}
\label{tab:notation}
\end{table}

% ======================================================================
\section{Joint Tissue--Phenotype Mean Matrix}
% ======================================================================

\label{sec:expansion-factors}

The one-step forward operator is assembled from three interpretable factors,
drawn per source state $z=(a,u)$, that separate the distinct biological processes
it bundles---expansion, tissue trafficking, and phenotype switching:
\begin{itemize}[leftmargin=1.4em, itemsep=2pt]
  \item $g_{a,u} \geq 0$: the per-source-state \textbf{expansion factor}---per-cell
        output over one step on the depth-normalized scale ($g>1$ expansion, $g<1$
        contraction / clonal loss); not an absolute or cross-tissue fold-change
        without an abundance anchor (Remark~\ref{rmk:anchor});
  \item $\pi_{(a,u)} \in \simplex{S-1}$: the destination-\textbf{tissue}
        distribution (\textbf{trafficking}), $\sum_{b\in\Ssp}\pi_{(a,u)}(b)=1$,
        with $\pi_{(a,u)}(a)$ the retained (stay-in-tissue) share;
  \item $\Phi_{(a,u),b}\in\simplex{K-1}$: the destination-\textbf{phenotype}
        distribution along the route to tissue $b$ (\textbf{switching}),
        $\sum_{v}\Phi_{(a,u),b}(v)=1$.
\end{itemize}

Assembling the factors gives the joint tissue--phenotype \emph{mean matrix}
$M \in \Reals_{\geq 0}^{L\times L}$ ($L = 39$), whose rows are source states and
whose columns are destination states:

\begin{keybox}[title={Factored mean matrix}]
\[
  M_{(a,u),(b,v)}
  \;=\;
  g_{a,u}\;\cdot\;\pi_{(a,u)}(b)\;\cdot\;\Phi_{(a,u),b}(v).
\]
\end{keybox}

For source state $z=(a,u)$ and destination $z'=(b,v)$, the entry $M_{zz'}$ is the
expected number of destination-$z'$ cells produced per source-$z$ cell over one
forward step,
\[
  M_{zz'}
  =
  \E[\,\text{cells in state } (b,v) \text{ at } r+1
        \text{ per cell in state } (a,u) \text{ at } r\,],
\]
and its rows are not constrained to sum to one---the row sum $g_z$ is freed to
carry expansion (a depth-normalized, not absolute cross-tissue, fold-change;
Remark~\ref{rmk:anchor}). Equivalently
$M=\operatorname{diag}(g)\,T$ with the row-stochastic
transition matrix
\[
  T_{(a,u),(b,v)} = \pi_{(a,u)}(b)\,\Phi_{(a,u),b}(v),
  \qquad
  g_z = \sum_{z'} M_{zz'} ,
\]
so the row sum recovers the expansion factor ($g_z = g_{a,u}$), $\pi_{(a,u)}$ is the
destination-tissue \emph{marginal} of the transition row, and $\Phi_{(a,u),b}$ its
destination-phenotype \emph{conditional} given the destination tissue.

On the interior where $g_z>0$ and $\pi_{(a,u)}(b)>0$ for every $b$, this is an
\emph{exact reparameterization}: each transition row equals its tissue marginal
times its phenotype conditional. Where a route carries no mass
($\pi_{(a,u)}(b)\to 0$), the corresponding switching factor $\Phi_{(a,u),b}$ is
not identified by the data and reverts to its prior
(Section~\ref{sec:posterior-sampling}). The construction serves to (i) place
biologically appropriate, \emph{separate} priors on expansion, trafficking, and
switching (Section~\ref{sec:priors})---in particular a persistence-favouring
trafficking prior---and (ii) expose these quantities as first-class parameters.
A conservative (non-expanding) model is the special case $g_z\equiv 1$.

\begin{remark}[Cross-tissue comparison of expansion]
\label{rmk:anchor}
The expansion factor $g_{a,u}$ is estimated on the depth-normalized scale---the
source enters as $\tilde{x}=x/\dsrc$ (each clone's cells divided by its tissue
sample's sequencing depth) and destination counts are scaled by $d^{\mathrm{dst}}$---so
$g$ measures per-cell output on a per-sample-depth basis, not an absolute cell-count
fold-change. Because those depths differ systematically across tissues, $g$ is
comparable across states within a source tissue but not across tissues: an apparent
cross-tissue difference in $g$ can reflect the tissues' differing depth regimes rather
than true proliferation. Recovering an absolute, cross-tissue fold-change needs an
\emph{abundance anchor} (per-sample total T-cell counts, cellularity, or a spike-in),
which we plan to add; like the missing-timepoint handling, this is a stated scope
limit, not a change to the model.
\end{remark}

\subsection{Destination intensity vector}
\label{sec:intensity}

For observation $j$ with depth-rescaled source $\tilde{x}_j$, the destination
\emph{intensity vector} is the pushforward of $\tilde{x}_j$ through $M$:

\begin{keybox}[title={Destination intensity}]
\[
  \mu_j = \tilde{x}_j M \in \Reals_{\geq 0}^{L},
  \qquad\text{equivalently}\qquad
  \mu_j(z')
  =
  \sum_{z\in\Zsp} \tilde{x}_j(z)\, M_{zz'} .
\]
\end{keybox}

The expected destination count in state $z'$ scales the intensity by the
destination depth of its tissue,
\[
  \E\!\left[ y_j(z') \mid M \right] = d_{j,s'}\, \mu_j(z'),
  \qquad s' = \mathrm{tissue}(z') .
\]
There is no simplex constraint on $\mu_j$ and no conditioning on a destination
total.

\subsection{Role of forward steps}

The same matrix $M$ is used for every patient $i$ and step $r \in \Rsp_i$. The
pair $(i,r)$ selects which source and destination vectors to build, not a
different matrix; all steps are repeated observations of a common one-step
forward process.

% ======================================================================
\section{Bayesian Treatment of the Mean Matrix}
% ======================================================================

The matrix $M$ is a latent non-negative kernel of per-cell destination
intensities over the joint tissue--phenotype state space.

\subsection{Forward generative reading}

Read forward, the additive mean structure follows from Poisson superposition.
Treating the depth-rescaled source as a population with $\tilde{x}_j(z)$ cells in
each state $z$, and letting each source cell in state $z$ contribute on average
$M_{zz'}$ cells to destination $z'$, the contributions sum across source states
and, scaled by the destination depth, give the expected destination count
\[
  \E\!\left[ y_j(z') \mid M \right]
  =
  d_{j,s'} \sum_{z\in\Zsp} \tilde{x}_j(z)\, M_{zz'}
  =
  d_{j,s'}\, \mu_j(z') ,
\]
the pushforward intensity of Section~\ref{sec:intensity}. This motivates the mean
structure; we do not take the Poisson as the emission law. Destination counts are
overdispersed relative to a Poisson---rate varies across clones and destination
states---so they are modelled with a Negative-Binomial whose mean is this
superposition mean and whose additional variance is a phenomenological
overdispersion term rather than a per-cell counting process
(Section~\ref{sec:likelihood}).

\subsection{Likelihood for destination counts}
\label{sec:likelihood}

\begin{keybox}[title={Observation likelihood}]
\[
  y_j(z') \;\sim\; \operatorname{NB2}\!\bigl( d_{j,s'}\, \mu_j(z'),\; \phi \bigr),
  \qquad s' = \mathrm{tissue}(z'),
\]
the Negative-Binomial in mean--concentration form: mean $d_{j,s'}\mu_j(z')$ and
variance $d_{j,s'}\mu_j(z') + \bigl(d_{j,s'}\mu_j(z')\bigr)^2 / \phi$, evaluated only
at destination states whose tissue is non-missing ($m_{j,s'} = 1$). The Poisson
superposition above is the overdispersion-free limit $\phi\to\infty$; the extra
term $\bigl(d_{j,s'}\mu_j(z')\bigr)^2/\phi$ absorbs overdispersion (rate variation)
beyond counting noise. The concentration $\phi=\exp(\log\phi)$ carries a log-scale
Normal prior and is a single global parameter, shared across all destination
states, patients, and steps. The latent matrix enters only through the intensity
$\mu_j$ (Section~\ref{sec:intensity}).
\end{keybox}

\subsection{Priors over the factors}
\label{sec:priors}

Priors are placed on the three factors of Section~\ref{sec:expansion-factors},
not on the $M$-entries, so each process carries a prior matched to its biology:
\[
  g_{a,u} \sim \operatorname{LogNormal}\!\bigl(\mu_g,\sigma_g^2\bigr),
  \qquad
  \pi_{(a,u)} \sim \Dir\!\bigl(\alpha_a\bigr),
  \qquad
  \Phi_{(a,u),b} \sim \Dir\!\bigl(\beta\bigr),
\]
each drawn independently over source states (and, for $\Phi$, destination
tissues), shared across all patients and steps; the Negative-Binomial
concentration has $\log\phi \sim \mathcal{N}(0,\sigma_\phi^2)$.
\begin{itemize}[leftmargin=1.4em, itemsep=2pt]
  \item \textbf{Expansion} $g_{a,u}$: a log-normal centred near unit expansion
        ($\mu_g\approx 0$) with a heavy right tail, so large expanders are not
        shrunk the way an exponential-tailed Gamma shrinks them.
  \item \textbf{Trafficking} $\pi_{(a,u)}$: a \emph{persistence-favouring}
        Dirichlet---the concentration vector $\alpha_a$ gives the source tissue
        ($b=a$) larger mass than the off-tissue destinations, so a priori a
        clone's descendants are expected to \emph{stay} rather than scatter. The
        concentration is indexed by source tissue $a$, so the stay mass can differ
        per tissue (e.g.\ a lower stay for CSF as a transit compartment).
  \item \textbf{Switching} $\Phi_{(a,u),b}$: a \emph{symmetric} Dirichlet over
        destination phenotypes ($\beta$ constant), with no a priori preference among
        fates---unlike trafficking, phenotype persistence is left to the data.
\end{itemize}

\begin{remark}[Why factored priors]
Placing priors on the factors states the model's biological expectations in the
same coordinates as its read-outs: the persistence-favouring $\Dir(\alpha_a)$ puts
prior mass on a clone's descendants remaining in the source tissue, and the
log-normal centres per-state expansion at unity ($g\approx 1$) with a heavy right
tail so large expanders are not shrunk.
\end{remark}

\begin{remark}[Support restrictions are an implementation detail]
The implementation may exclude biologically implausible transitions by removing
the corresponding destinations from the relevant $\pi$ or $\Phi$ support. This
sparsifies $M$ without changing the form of the model; sums over destination
states below are read over the retained support. The likelihood-support guard is
unchanged: every non-missing destination $z'$ with $y_j(z')>0$ must be reachable
($M_{zz'}>0$) from some source state in the support of $\tilde{x}_j$.
\end{remark}

\subsection{Posterior predictive destination counts}

Because $M$ is latent, the intensity $\mu_j = \tilde{x}_j M$ is also uncertain.
Posterior draws of $M$ induce replicated destination counts
\[
  y_j^{\mathrm{rep}}(z') \;\sim\; \operatorname{NB2}\!\bigl( d_{j,s'}\, \mu_j(z'),\; \phi \bigr)
\]
over destination states with $m_{j,s'} = 1$.

\subsection{Backward attribution}
\label{sec:backward}

The reverse question---which source state produced an observed destination
cell---is a posterior attribution. For observation $j$,
\[
  w_{jzz'}
  =
  \frac{\tilde{x}_j(z)\, M_{zz'}}
       {\sum_{\tilde z\in\Zsp} \tilde{x}_j(\tilde z)\, M_{\tilde z z'}}
  =
  \frac{\tilde{x}_j(z)\, M_{zz'}}{\mu_j(z')} ,
\]
the same form as the forward intensity with source and destination roles
reversed. Being a function of $M$, it is a posterior read-out: evaluated at the
fitted $\widehat{M}$, or per draw to carry uncertainty, it attributes observed
destination cells back to their likely sources.

% ======================================================================
\section{Full Generative Story}
\label{sec:generative}
% ======================================================================

The full model is a model for the one-step expansion and redistribution of clone
mass across tissue--phenotype states. It can be summarized as follows.

\begin{keybox}[title={Generative model}]
The state space $\Zsp$ ($|\Zsp| = L$) and the fixed inputs---each observation's
depth-rescaled source $\tilde{x}_j$, the destination depths $d_{j,s'}$, and the
missingness flags $m_{j,s'}$---are given. The model generates the destination
counts as follows.
\begin{enumerate}[leftmargin=2em, itemsep=4pt]
  \item \textbf{Operator factors.} For each source state $z=(a,u)$, draw
        independently the expansion, trafficking, and switching factors
        \[
          g_{a,u} \sim \operatorname{LogNormal}(\mu_g,\sigma_g^2),
          \quad
          \pi_{(a,u)} \sim \Dir(\alpha_a),
          \quad
          \Phi_{(a,u),b} \sim \Dir(\beta)\ \ (b\in\Ssp),
        \]
        and the dispersion $\log\phi \sim \mathcal{N}(0,\sigma_\phi^2)$; assemble the
        mean matrix
        $M_{(a,u),(b,v)} = g_{a,u}\,\pi_{(a,u)}(b)\,\Phi_{(a,u),b}(v)$, shared
        across all patients and steps.

  \item \textbf{Observation.} For each observation $j = (i,r,c)$, push the source
        forward to the intensity $\mu_j = \tilde{x}_j M$, and for each destination
        state $z'$ with $m_{j,s'} = 1$, draw
        \[
          y_j(z') \;\sim\; \operatorname{NB2}\!\bigl( d_{j,s'}\, \mu_j(z'),\; \phi \bigr),
          \qquad s' = \mathrm{tissue}(z') .
        \]
\end{enumerate}
\end{keybox}

The dependency structure is shown in Figure~\ref{fig:plate}.

\begin{figure}[htbp]
  \centering
  \includegraphics[width=0.78\linewidth]{plate_joint_transition.pdf}

  \smallskip
  {\small
  \begin{tabular}{@{}ll@{}}
    \toprule
    \textbf{Node} & \textbf{Meaning}\\
    \midrule
    $x_{irc}$ & observed source counts at the start of the step (over the $L$ states)\\
    $\tilde{x}_{irc}=x_{irc}/\dsrc_{ir,s}$ & depth-rescaled source (deterministic; folded into the $x\!\to\!y$ edge)\\
    $y_{irc}$ & observed destination counts in tissue $s$ at the end of the step\\
    $g_{z}$ & latent expansion factor per source state (log-normal)\\
    $\pi_{(a,u)}$ & latent destination-tissue distribution --- trafficking (Dirichlet)\\
    $\Phi_{(a,u),b}$ & latent destination-phenotype distribution per route --- switching (Dirichlet)\\
    $M_{(a,u),(b,v)}=g\,\pi\,\Phi$ & assembled mean matrix, shared across all patients and steps\\
    $\phi$ & Negative-Binomial concentration (log-Normal)\\
    $d_{irs}$ & destination sequencing depth, exposure (fixed)\\
    $(\mu_g,\sigma_g,\alpha_a,\beta,\sigma_\phi)$ & prior hyperparameters (fixed)\\
    \bottomrule
  \end{tabular}}

  \caption{Plate diagram of the generative model (node key above). The operator
    factors---expansion $g_z$, trafficking $\pi_{(a,u)}$, and switching
    $\Phi_{(a,u),b}$---are drawn once per source state from their priors and
    assembled into
    $M_{(a,u),(b,v)}=g_{a,u}\,\pi_{(a,u)}(b)\,\Phi_{(a,u),b}(v)$; this block sits
    \emph{outside} the patient plate, so a single $M$ (and dispersion $\phi$) is
    shared across all patients and steps. The observation side nests patient $i$,
    step $r\in\Rsp_i$, clone $c\in\Csp_{ir}$, and tissue $s\in\Ssp$. The
    depth-rescaled source $\tilde{x}_{irc}$ is pushed through $M$ to the intensity
    $\mu_{irc}=\tilde{x}_{irc}M$, and each destination count $y_{irc}$ is a
    Negative-Binomial draw with mean $d_{irs}\,\mu_{irc}$ and concentration $\phi$,
    evaluated only at non-missing tissues (masking is a support restriction, not a
    drawn node). Double circles ($x_{irc}$, $y_{irc}$) mark observed quantities;
    solid dots (prior hyperparameters, $d_{irs}$) mark fixed inputs; the
    depth-rescaling $\tilde{x}_{irc}=\mathrm{Rescale}(x_{irc})$ is folded into the
    $x\!\to\!y$ edge.}
  \label{fig:plate}
\end{figure}

% ======================================================================
\section{Inference}
% ======================================================================

Inference estimates the posterior over the shared mean matrix $M$. The
observed data are the clone-level forward observations $\Dsp$; the depth-rescaled
sources $\tilde{x}_j$, the depths $d_{j,s'}$, and the missingness flags
$m_{j,s'}$ are fixed inputs, as are the prior hyperparameters
$(\mu_g,\sigma_g,\alpha_a,\beta,\sigma_\phi)$.

\subsection{Learned quantity}

The main learned object is the posterior $p(M \mid \Dsp)$. A single-entry
summary is
\[
  \E\!\left[ M_{(a,u),(b,v)} \mid \Dsp \right],
\]
the posterior mean number of destination-$(b,v)$ cells produced per
source-$(a,u)$ cell over one forward step.

\subsection{Posterior sampling}
\label{sec:posterior-sampling}

The posterior is not available in closed form. The unknowns are the operator
factors $\{g_{a,u},\,\pi_{(a,u)},\,\Phi_{(a,u),b}\}$ and the dispersion $\log\phi$;
the mean matrix $M_{(a,u),(b,v)}=g_{a,u}\pi_{(a,u)}(b)\Phi_{(a,u),b}(v)$ and the
intensities $\mu_j=\tilde{x}_j M$ are deterministic functions of them, and the
mean--concentration Negative-Binomial is written directly on the aggregated
destination counts, so no auxiliary allocation variables are introduced. We draw
$\{g,\pi,\Phi,\log\phi\}^{(1)},\ldots,\{g,\pi,\Phi,\log\phi\}^{(N)}$ by gradient-based
Markov chain Monte Carlo (NUTS) in the unconstrained reparameterization---positive
$g$ through a log transform, the simplices $\pi,\Phi$ through their standard
bijections. The log-posterior is smooth: the intensity is linear in $M$ and $M$ is
a smooth product of the factors, and the log-normal, Dirichlet, and
Negative-Binomial densities are smooth on their supports. Each $\Phi_{(a,u),b}$
enters the likelihood only through the trafficking mass $\pi_{(a,u)}(b)$ routed to
tissue $b$, so off-route switching factors (small $\pi_{(a,u)}(b)$) are only weakly
informed and prior-dominated; the proper Dirichlet priors keep these directions
well-behaved.

\subsection{Estimated mean matrix and uncertainty}

The fitted mean matrix is the posterior mean $\widehat{M} = \E[M \mid \Dsp]$,
estimated by the sample average $\tfrac{1}{N}\sum_{n} M^{(n)}$ over the $N$ draws; its row sums are
free---$\widehat{M}$ is not normalized across destinations. Credible intervals
for any entry follow from the sample quantiles. Every read-out below is a
deterministic function of $M$, so evaluating it on each draw $M^{(n)}$ yields its
full posterior, summarized by a posterior mean and credible interval.

% ======================================================================
\section{Posterior Summaries}
% ======================================================================

The fitted $\widehat{M}$ yields, for any source tissue and phenotype, three
per-step quantities---where a clone's mass traffics across tissues, whether it
expands or contracts, and how its phenotype mix reshapes---each a dimensionless
fraction or factor over one step, not a per-time rate.

Define $\rho^a(u)$ as the empirical source-cell phenotype distribution in tissue
$a$, pooled across clones and steps; it collapses source phenotypes in the
tissue-level rows below.

\begin{table}[h]
\centering
\small
\renewcommand{\arraystretch}{1.4}
\begin{tabular}{@{}p{2.15cm}p{2.35cm}p{3.2cm}p{5.3cm}@{}}
\toprule
\textbf{Quantity} & \textbf{Grain} & \textbf{Reads as} & \textbf{Expression}\\
\midrule
Traffic (resolved)
  & (tissue, phenotype) $\to$ tissue
  & fraction of mass, per step
  & $\pi_{(a,u)}(b)$\\
Traffic (collapsed)
  & tissue $\to$ tissue
  & fraction of mass, per step
  & $\sum_{u}\rho^a(u)\,\pi_{(a,u)}(b)$\\
Traffic rate
  & tissue
  & retention vs.\ out, per step
  & stay $=\pi_{(a,u)}(a)$, out $=\sum_{b\neq a}\pi_{(a,u)}(b)$; sum to $1$\\
Expansion factor
  & (tissue, phenotype)
  & per-cell output, per step ($>1$ expand, $<1$ contract; Rmk.~\ref{rmk:anchor})
  & $g_{a,u}=\sum_{z'}\widehat{M}_{(a,u),z'}$\\
Clone expansion
  & clone
  & clone-mean over source states
  & $\sum_{z}\theta(z)\,g_z$\\
Switching (resolved)
  & (tissue, phenotype) $\to$ tissue
  & phenotype distribution, per step
  & $\Phi_{(a,u),b}(v)$\\
Switching (collapsed)
  & tissue $\to$ tissue
  & phenotype distribution, per step
  & $\sum_{u}\rho^a(u)\,\Phi_{(a,u),b}(v)$\\
Switch rate
  & tissue route
  & switched vs.\ kept, per step
  & switched $=\sum_{v\neq u}\Phi_{(a,u),b}(v)$, kept $=\Phi_{(a,u),b}(u)$\\
\bottomrule
\end{tabular}
\caption{Per-step read-outs of the fitted mean matrix $\widehat{M}$.}
\label{tab:readouts}
\end{table}

The resolved read-outs are the model's own factors---trafficking $\pi_{(a,u)}$,
switching $\Phi_{(a,u),b}$, and the expansion factor $g_{a,u}$
(Section~\ref{sec:expansion-factors})---read \emph{directly}, each with its own
posterior mean and credible interval; the collapsed read-outs are their
$\rho^a$-weighted tissue aggregates. $\theta(z)$ is a clone's source
composition---its source cells as a distribution over states.

\end{document}
